\subsection{Uniform first-passage fibres and separated endpoint intervals}
\label{subsec:Allikvere-uniform-fibres}

Allikvere's Section~6 \cite{Allikvere2026v2} replaces the logarithmic
initial distribution by the uniform distribution on odd integers in
$[y,y^\alpha]$. The upper end of this interval carries almost all the
uniform mass. The time interval must consequently retain that end:
\[
 L=\log x,\quad d=\log(4/3),\quad \alpha=1.001,\quad
 m_0=\lfloor L/100000\rfloor,\qquad
 \widetilde I_y=
 [\log(y/x)/d+L^{4/5},\ \log(y^\alpha/x)/d+L^{4/5}].
\]
Here $y=x^\alpha$ or $x^{\alpha^2}$. Let $E'$ be the full set of
positive odd $M$ which first enter $[1,x]$ after exactly $m_0$ Syracuse
steps and satisfy
$M_-=xe^{dm_0-L^{7/10}}\le M\le M_+=xe^{dm_0+L^{7/10}}$.
The word ``full'' matters: an arbitrary restriction on the final value
can remove the interval structure within a valuation cylinder. Absolute
errors will be bounded on the full set before taking any such subset.

Set $\mathcal I_y=(\widetilde I_y-m_0)\cap\mathbb Z$ and
$R=\max\mathcal I_y$. Then $r\asymp_\alpha L$ on $\mathcal I_y$,
$3^R\le x^{0.0116}$, and $M_->4\,3^{2R}$ for sufficiently large $x$.
Indeed $(\alpha^3-1)\log3/d<0.0116$, whereas
$\log M_-=L+dm_0-L^{7/10}$.

\begin{lemma}[integer separation of endpoint incidences]
\label{lem:Allikvere-integer-endpoint}
For $R\ge1$, $M>4\,3^{2R}$ and $v>0$, at most one pair
$(r,s)\in\{1,\ldots,R\}\times\mathbb Z$ satisfies
$M-3^r\le v3^r2^{-s}\le M$.
\end{lemma}
\begin{proof}
Put $Q=3^R$. All candidate values lie in $[M-Q,M]$, an interval of
endpoint ratio less than $2$, so two different $s$ cannot occur at the
same $r$. If $r_2>r_1$, write $z_i=v3^{r_i}2^{-s_i}$,
$q=r_2-r_1\ge1$, $p=s_2-s_1$. From $z_2/z_1=3^q/2^p<2$ we have
$p\ge1$. The unequal positive integers $3^q$ and $2^p$ would satisfy
\[
 1\le|3^q-2^p|=2^p|z_2-z_1|/z_1
 \le Q^2M/(M-Q)^2\le4Q^2/M<1.
\]
Here $2^p=3^qz_1/z_2\le QM/(M-Q)$, $|z_2-z_1|\le Q$ and
$M-Q>M/2$. This contradiction proves the lemma.
\end{proof}

The complete finite interval geometry is the map
\[
 (r,s)\longmapsto[(M-3^r)2^s/3^r,\ M2^s/3^r]
\]
in the endpoint variable $v$: these closed intervals are pairwise
disjoint. Both endpoints and the full allowed offset $3^r$ are retained.

For a positive word $\mathbf a$ of length $r$, total $s$, and forward
offset $F_r(\mathbf a)$, the input ending at $M$ is exactly
$N=2^s(M-F_r(\mathbf a))/3^r$ when the integrality congruence holds.
The comparison point is $N^0=2^sM/3^r$, with
$0<F_r(\mathbf a)\le3^r$. If the closed true window
$y\le N\le y^\alpha$ and half-open replacement $y<N^0\le y^\alpha$
disagree, at least one of $v=y,y^\alpha$ lies between them. Such a pair
therefore satisfies the endpoint lemma. This bounds the number of
possible $(r,s)$ per endpoint and per $M$ over the entire time interval;
no linear-form-in-logarithms estimate is needed here.

\begin{proposition}[absolute window-replacement bound]
\label{prop:Allikvere-window-replacement}
After multiplication by $2/y^\alpha$, the number of all
$(r,\mathbf a,M)$ with $r\in\mathcal I_y$, $M\in E'$,
$F_r(\mathbf a)\equiv M\pmod{3^r}$ and differing true/replaced indicators
is $O_\alpha(L^{-1/2})$.
\end{proposition}
\begin{proof}
The full return band is a union of $P_x\ll x^{10^{-4}}$ interval pieces
in distinct dyadic valuation cylinders. To count the pieces, the
pre-crossing inequality and $F_{m_0-1}<x/2$ give
$b_1+\cdots+b_{m_0-1}\le2m_0+O(L^{7/10})$, so there are at most
$4^{m_0}e^{O(L^{7/10})}$ initial words. The final valuation is $O(L)$
by $2^{b_{m_0}}\le3^{m_0+1}(M+1)$. This gives the stated $P_x$.
The exact odd cylinder modulus for a word of total $B$ is $2^{B+1}$.
Consequently its count in any interval and any residue modulo $3^r$
differs from $3^{-r}$ times its total count by at most $1$, by coprimality.
The full band's discrepancy is therefore at most $P_x$.

For fixed $(r,s)$ the differing pairs lie in the union of the two
intervals $[v3^r2^{-s},v3^r2^{-s}+3^r]$ in $M$.
Applying the preceding discrepancy to its at most two components gives
the uniform part of the incidence count and an error $O(P_x)$ per residue.
The conditioned residue histogram of all words has total mass $1$.
Use the identity
$\binom{s-1}{r-1}=3^r2^{s-r\log_2 3}p_r(s)$, where
$p_r(s)=\binom{s-1}{r-1}2^{-s}$.
For a contributing pair $2^{s-r\log_2 3}\le2y^\alpha/M_-$;
thus the discrepancy part is at most
$CP_x3^R|\mathcal I_y|/M_-\ll x^{-0.98}$ after summing $p_r(s)$ first.

For the uniform part interchange the nonnegative sums and fix $M,v$.
There is at most one $(r,s)$ by the endpoint lemma, with
$2^{s-r\log_2 3}\le2v/M$ and $p_r(s)\ll r^{-1/2}\ll_\alpha L^{-1/2}$.
The latter bound follows by the ratio
$p_r(s+1)/p_r(s)=s/[2(s-r+1)]$ and Stirling at the central mode.
Summing the two endpoints gives
$CL^{-1/2}(1+y^{1-\alpha})\sum_{M\in E'}1/M\ll L^{-1/2}$,
by the previously proved uniform harmonic coefficient bound.
\end{proof}

The source's \texttt{lem:winrep} uses Rhin's estimate and obtains
$L^{-1/2+o(1)}$. The proved improvement removes that dependency from
endpoint replacement, not from the separate averaged kernel asymptotic.

For precision, the remaining error transfer has also been written out in
\path{research_program/natural_density_audit_20260928/audit.tex},
Propositions \texttt{prop:kernel-upper} and \texttt{prop:fibre-master}.
For $u=\log_2(y^\alpha/M)$, $u_0=\log_2(y/M)$, define
\[
 D_y(M)=\sum_{r\in\mathcal I_y}3^{-r}
       \sum_{r\log_2 3+u_0<s\le r\log_2 3+u}\binom{s-1}{r-1}.
\]
The exact hockey-stick identity and the uniform Stirling bound
$p_r(s)\le Cr^{-1/2}e^{-c(s-2r)^2/r}$ for $s/r$ in the relevant compact
interval prove $D_y(M)\ll y^\alpha/M$ independently of equidistribution.
Specifically, the upper-truncated sum is at most
$2^u(s_r/r)p_r(s_r)$, $s_r=\lfloor r\log_2 3+u\rfloor$;
the Gaussian sum in $r$ has width $O(\sqrt L)$ and height $O(L^{-1/2})$.

Let $\mathcal A_r$ retain all prefix inequalities
$|a_1+\cdots+a_j-2j|<L^{3/5}$. With uniform odd input $\widetilde N_y$,
$m_1=\lfloor L^{2/5}\rfloor$ and $\rho_{m_1}=\operatorname{Law}(X_{m_1})$,
the proved fibre formula is
\[
 \sum_{r\in\mathcal I_y}\sum_{\mathbf a\in\mathcal A_r}
 \mathbb P(\operatorname{Aff}_{\mathbf a}(\widetilde N_y)=M)
 =\frac{2}{y^\alpha}3^{m_1}\rho_{m_1}(M\bmod3^{m_1})D_y(M)+R_y(M),
 \qquad\sum_{M\in E'}|R_y(M)|\ll L^{-1/4}.
\]
The proof pairs the conditioned mixing error with actual band populations,
using residue discrepancy and the kernel upper bound. Exceptional words
and noncentral totals are controlled with the same population weights;
the restriction $\mathcal A_r$ is not asserted invariant under reversal.
The leading constant is evaluated next. The actual-input reduction and
natural-density/time assembly remain separate from these component estimates.

\subsubsection{The phase average without freezing the weight}
\label{subsec:Allikvere-phase-average}

Put $\lambda=\log_2 3$, $a=2-\lambda$, $\mu=8.616$ and
$\beta=1-1/\mu$. Rhin's original proposition (8), p.~160
\cite{Rhin1987}, implies $\|q\lambda\|\ge c_\lambda q^{-(\mu-1)}$:
take the nearest integer $p$ to $q\lambda$, use $p\le2q$, and absorb
the finitely many small denominators. The form cannot vanish since
$3^q\ne2^p$. The original auxiliary-polynomial computation is used
as a published input, not claimed as independently replayed.

Erd\H{o}s--Tur\'an and Koksma, in the precise forms of
\cite[Ch.~2, Theorems~2.5 and 5.1]{KN1974}, give the centered sum of
a real BV circle function $f$ on any $k$ consecutive rotation points
the bound $C_\lambda\operatorname{Var}(f)k^\beta$.
Indeed the geometric sum at frequency $h$ is at most
$(2\|h\lambda\|)^{-1}$, so unnormalized discrepancy is at most
$C_\lambda(k/H+H^{\mu-1})$; choose $H=\max(1,\lfloor k^{1/\mu}\rfloor)$.
For a weight on $[A,B]\cap\mathbb Z$, set
$S_j=\sum_{n=A}^j(f(\{n\lambda+u\})-\int f)$.
The exact identity
\[
 \sum_{n=A}^B w(n)(S_n-S_{n-1})
 =w(B)S_B+\sum_{n=A}^{B-1}(w(n)-w(n+1))S_n
\]
therefore bounds the weighted error by
$C_\lambda\operatorname{Var}(f)(B-A+1)^\beta
\sum_{n\in\mathbb Z}|w(n+1)-w(n)|$,
where $w$ is extended by zero. Both endpoint jumps are retained.
This removes the block-freezing error from the source's \texttt{lem:E}.
There is no deterministic phase-weight independence requirement.

For $z=u/a\asymp_\alpha L$ introduce
\[
 g_z(t)=(\pi t)^{-1/2}
       e^{-a^2(t-z)^2/(4t)}\quad(t>0),\qquad g_z(0)=0.
\]
Its integral is exactly $2/a$. In fact $t=s^2$, followed by the
bijection $v=s-z/s$, has inverse
$s=(v+\sqrt{v^2+4z})/2$ and derivative
$(1+v/\sqrt{v^2+4z})/2$. Consequently
\[
 \int_0^\infty g_z(t)dt
 =\pi^{-1/2}\int_{\mathbb R}e^{-a^2v^2/4}
           (1+v/\sqrt{v^2+4z})dv=2/a.
\]
The second term is odd and absolutely integrable.
Also $4t^2(\log g_z)'=a^2z^2-2t-a^2t^2$ has one positive zero,
$\sqrt{z^2+a^{-4}}-a^{-2}\ge z/2$ for large $z$.
Thus $g_z$ is unimodal with maximum at most $\sqrt{2/(\pi z)}$.
The zero-extended samples on
$J_z=\{r:|r-z|\le\sqrt z\log L\}$ have variation twice their
maximum, hence $O_\alpha(L^{-1/2})$. Comparing each unit interval
with its left endpoint bounds the full sum-integral error by
$\int|g'_z|=2\max g_z$. The omitted tails are $O_A(L^{-A})$:
use a Gaussian bound when $z/2\le r\le2z$, and exponents
$-cz^2/r$ and $-cr$ on the two outer regions. It follows that
$\sum_{J_z}g_z(r)=2/a+O_\alpha(L^{-1/2})$.

\begin{proposition}[leading kernel with a stronger component rate]
\label{prop:Allikvere-leading-kernel}
Uniformly for every real $M\in[M_-,M_+]$ and both specified $y$,
\[
 D_y(M)=\frac{y^\alpha}{M\log(4/3)}
          (1+O_{\alpha,c}(L^{-c}))
 \qquad(0<c<1/17.232).
\]
In particular $c=1/18$ is admissible.
\end{proposition}
\begin{proof}
The exact two endpoints of the inner sum are
$s_r=\lfloor r\lambda+u\rfloor$ and
$k_r=\lfloor r\lambda+u_0\rfloor$.
Hockey-stick summation gives
$3^{-r}(\binom{s_r}{r}-\binom{k_r}{r})$.
The bottom sum divided by $2^u$ is
$O_\alpha(L2^{-(u-u_0)})$, exponentially small in $L$.
The top term divided by $2^u$ is
$2^{-\theta_r}(s_r/r)p_r(s_r)$, $\theta_r=\{r\lambda+u\}$.

Write $M=xe^{dm_0+\sigma}$, $|\sigma|\le L^{7/10}$. Then
$z=\log(y^\alpha/x)/d-m_0-\sigma/d$.
The top margin in $\widetilde I_y-m_0$ is
$L^{4/5}+\sigma/d\ge L^{4/5}/2$; the bottom margin is
$(\alpha-1)\log y/d-L^{4/5}-\sigma/d\gg_\alpha L$.
Hence $J_z\subset\mathcal I_y$. The Gaussian bound already proved
discards its complement with error $O_A(L^{-A})$ after division by $2^u$.

On $J_z$, Stirling and Taylor expansion of
$H(b)=b\log b-(b-1)\log(b-1)-b\log2$ at $b=2$ give
\[
 p_r(2r+\delta)=\frac{e^{-\delta^2/(4r)}}{\sqrt{4\pi r}}
 (1+O(r^{-1}+|\delta|/r+|\delta|^3/r^2)).
\]
Here $H''(2)=-1/2$ and the factorial prefactor is
$(2\pi r b(b-1))^{-1/2}$.
Since $\delta=-a(r-z)-\theta_r$ and $0\le\theta_r<1$,
$(s_r/r)p_r(s_r)=g_z(r)(1+O_\alpha(L^{-1/2}(\log L)^3))$.
The bounded mass of $g_z$ makes this an absolute summed error.
Apply the preceding Abel estimate to $f(t)=2^{-t}$ on the circle:
its variation is $1$ with the endpoint jump, and its integral is
$1/(2\log2)$. The support has $O_\alpha(\sqrt L\log L)$ points.
Thus
\[
 2^{-u}D_y(M)=\frac1{a\log2}
 +O_\alpha(L^{-1/(2\mu)}(\log L)^{1-1/\mu}
                   +L^{-1/2}(\log L)^3).
\]
Use $a\log2=d$ and $2\mu=17.232$ to conclude.
\end{proof}

Substitution into the actual fibre formula gives
\[
 \sum_{M\in E'}\left|W_y(M)
  -\frac{2\,3^{m_1}}{d}\frac{\rho_{m_1}(M)}M\right|
  \ll_{\alpha,c}L^{-c}\quad(0<c<1/17.232).
\]
The proof uses the already established weighted harmonic bound to sum
the pointwise kernel errors. The same bound survives any pushforward
of these finite measures, by the triangle inequality on each fibre.
The following argument identifies these measures with the first-passage
probability and carries the rate to that law. The global natural-density
iteration and timed endpoint are separate subsequent arguments.

\paragraph{Actual uniform input and the probability law.}
\label{par:Allikvere-actual-input}
Let $n_0=\lfloor L/(10\log2)\rfloor$ and let $\widetilde N_y$
be uniform among the positive odd integers in $[y,y^\alpha]$.
Write $\operatorname{Pass}_x(N)=\operatorname{Syr}^{T_x(N)}(N)$
when $T_x(N)<\infty$, and set its value to $1$ otherwise.

\begin{proposition}[uniform first-passage stabilisation with the stronger rate]
\label{prop:Allikvere-uniform-stabilisation}
Let
\[
 \Pi_x(M)=\frac{2\,3^{m_1}}d\frac{\rho_{m_1}(M\bmod3^{m_1})}{M},
 \qquad q_x:E'\longrightarrow(2\mathbb N+1)\cap[1,x],
 \quad q_x(M)=\operatorname{Syr}^{m_0}(M).
\]
For every $0<c<1/17.232$ and sufficiently large $x$,
$\Pi_x(E')=1+O_{\alpha,c}(L^{-c})>0$. The probability
$Q_x=(q_x)_*\Pi_x/\Pi_x(E')$ is independent of $y$, and
\[
 \|\operatorname{Law}(\operatorname{Pass}_x(\widetilde N_y))-Q_x\|_1
 \ll_{\alpha,c}L^{-c}
 \quad(y=x^\alpha,\ x^{\alpha^2}).
\]
Thus the two laws differ by $O_{\alpha,c}(L^{-c})$ in $\ell^1$.
Also $\mathbb P(T_x(\widetilde N_y)=\infty)\ll x^{-\eta}$ for some
$\eta>0$. The norm is the sum of absolute atom masses.
\end{proposition}
\begin{proof}
Put $q=3n_0$ and $K_y=\#((2\mathbb N+1)\cap[y,y^\alpha])$.
Counting each odd residue modulo $2^q$ gives total discrepancy
$O(2^q/K_y)\ll2^{-q}$, since $4^{3n_0}\le x^{3/5}\le y$.
Tao's Proposition~1.9 \cite{Tao2026v7} applies with $c_0=1$:
the length-$n_0$ valuation vector is $O(x^{-\eta})$ close in
$\ell^1$ to independent geometric variables of masses $2^{-k}$.
One can also see the bound directly. A word of total $s<q$ is
one odd class modulo $2^{s+1}$ and has probability
$2^{-s}+O(2^{-q})$; there are $\binom{q-1}{n_0}$ such words.
Stop a higher-total word just before its total reaches $q$.
A length-$k<n_0$ prefix of total less than $q$ has
$\binom{q-1}{k}$ possibilities, and divisibility of the next affine
numerator fixes one class modulo $2^q$. Its probability is
$O(2^{-q})$. The binomial lower tail therefore bounds both the
accumulated low-word error and the high-total mass by $e^{-cn_0}$.

For a centered geometric variable the moment-generating function is
$e^{-t}/(2-e^t)$; exponential Markov gives
$\mathbb P(|S_j-2j|\ge h)\le2e^{-c\min(h^2/j,h)}$.
Let $G$ be the actual event that every prefix through $n_0$ has
$|S_j-2j|<L^{3/5}$. The preceding law comparison and union bound
give $\mathbb P(G^c)\ll e^{-\kappa L^{1/5}}$.
The exact maximal-moment calculation in
Corollary~\ref{cor:Mazur-allikvere-prefix} strengthens this to
$(2+o(1))\exp[-(5\log2/2)L^{1/5}]$, with both closed input
endpoints retained.
On $G$ the exact affine expression yields
\[
 \log\operatorname{Syr}^j(N)
 =\log N-dj-(\log2)(S_j-2j)+O(x^{-3/4}).
\]
Indeed its linear part exceeds $x^{0.95}$ and its offset is at most
$3^{n_0}<x^{0.16}$ for large $x$; the inequalities
$3^2>2^3$ and $3^5<2^8$ supply the strict exponent margins.
The orbit crosses by $n_0$ because
$\alpha^3-d/(10\log2)<\alpha^3-1/25<1$.
Applying the drift formula on the two sides of the crossing gives
$T_x(N)=\log(N/x)/d+O_\alpha(L^{3/5})$.
The upper endpoint of $\widetilde I_y$ is respected on $G$.
Failure of the lower endpoint forces
$N\le y e^{2dL^{4/5}}$, whose uniform probability is
$O(y^{1-\alpha}e^{2dL^{4/5}}+1/K_y)=O(x^{-\eta})$.
Hence $\mathbb P(T_x\notin\widetilde I_y)\ll e^{-\kappa L^{1/5}}$.

For each actual input define the counting-weight measure
\[
 H_N=\sum_{n\in\widetilde I_y\cap\mathbb Z}
 {\bf1}_{\vec a^{(n-m_0)}(N)\in\mathcal A_{n-m_0}}\,
 {\bf1}_{\operatorname{Syr}^{n-m_0}(N)\in E'}\,
 \delta_{q_x(\operatorname{Syr}^{n-m_0}(N))}.
\]
The exact inverse fibres, retaining the true input interval, give
$\mathbb EH_N=(q_x)_*W_y$.
For $N\in G$ and $n\in\widetilde I_y$, put $r=n-m_0$.
If $T_x(N)=n$, drift puts $M=\operatorname{Syr}^r(N)$ inside
$[M_-,M_+]$ and $T_x(M)=m_0$. Conversely, if $M\in E'$, then
for every $j<r$,
\[
 \log\operatorname{Syr}^j(N)
 \ge\log M+d(r-j)-C_\alpha L^{3/5}
 \ge\log x+dm_0-L^{7/10}-C_\alpha L^{3/5}>\log x.
\]
The definition of $E'$ controls indices $r,\ldots,n$, so the first
crossing is precisely $n$. Consequently $H_N$ equals the passage
atom on $G\cap\{T_x\in\widetilde I_y\}$ and is zero on the rest of $G$.
It always has at most $J=\#(\widetilde I_y\cap\mathbb Z)=O_\alpha(L)$
atoms counted with multiplicity. Thus
\[
 \|\operatorname{Law}(\operatorname{Pass}_x(\widetilde N_y))
              -(q_x)_*W_y\|_1
 \le\mathbb P(T_x\notin\widetilde I_y,G)+(J+1)\mathbb P(G^c)
 \ll L e^{-\kappa L^{1/5}}\ll_B L^{-B}
\]
for every fixed $B>0$.
The proved fibre-profile estimate and contraction of pushforward
now show that the law is $O(L^{-c})$ from $(q_x)_*\Pi_x$.
Taking total masses proves $\Pi_x(E')=1+O(L^{-c})>0$.
Dividing by this proved positive mass changes the measure in
$\ell^1$ by exactly $|\Pi_x(E')-1|$, proving the asserted comparisons.

Finally $\mathbb P(S_{n_0}\le1.9n_0)\ll x^{-\eta}$ by the same
geometric exponential bound and actual-law comparison.
Outside this event the affine bound gives
$\operatorname{Syr}^{n_0}(N)\le
C x^{\alpha^3+(\log_2 3)/10-0.19}+x^{(\log_2 3)/10}<x$,
since $\alpha^3-0.03<1$ and $\log_2 3<1.6$.
Thus $T_x=\infty$ has the stated power bound.
\end{proof}

This proves the uniform first-passage statement in Allikvere's
source label \texttt{thm:D}, using his extended window and measure
change, Tao's valuation and deterministic passage mechanisms, and
the stronger kernel estimate above. The subsequent sections carry
the rate to timed fixed-target counts and natural density.
Full source locators are in the standalone audit, claims ND-020--ND-027.

\subsection{The scale iteration and the fixed-target estimate}

Retain $O_z=(2\mathbb N+1)\cap[1,z]$. Write $U_y$ for the uniform probability on the positive odd integers
in $[y,y^\alpha]$, whenever this set is nonempty. All probabilistic
applications below have $y$ larger than a fixed constant. The final
counting statements hold for every real cutoff $X\ge2$; no uniform
probability on an empty short interval is needed. Fix
$0<c<1/17.232$, and retain $\alpha=1001/1000$ and $d=\log(4/3)$.

\begin{lemma}[Nested passage, with failure kept separate]
\label{lem:Allikvere-nested-passage}
For $1\le z\le w$ and a positive odd integer $N$ with $T_z(N)<\infty$,
put $M=\operatorname{Pass}_w(N)$. Then
\[
 T_w(N)\le T_z(N),\qquad
 T_z(N)=T_w(N)+T_z(M),\qquad
 \operatorname{Pass}_z(N)=\operatorname{Pass}_z(M).
\]
For any subset $E$ of the positive odd integers,
\[
 \mathbb P(\operatorname{Pass}_z(N)\in E)
 -\mathbb P(T_z(N)<\infty,\operatorname{Pass}_z(N)\in E)
 =\mathbf1_{1\in E}\mathbb P(T_z(N)=\infty).
\]
The second identity holds for any random positive odd input.
\end{lemma}
\begin{proof}
Before the first entry below $w$, every iterate exceeds $w\ge z$.
Starting at that entry, the first entry below $z$ is exactly the
remaining part of the same orbit. This proves the three deterministic
identities, including the case $M\le z$, when the remaining time is zero.
On failed passage the defined value of $\operatorname{Pass}_z$ is $1$;
splitting the probability according to finite or infinite passage proves
the last equality.
\end{proof}

\begin{proposition}[The error survives the scale ladder]
\label{prop:Allikvere-untimed-ladder}
For every integer $N_0\ge2$ and every sufficiently large $y$,
\[
 U_y\{N:\operatorname{Syr}^k(N)>N_0\text{ for all }k\ge0\}
 \ll_c(\log N_0)^{-c}.
\]
\end{proposition}
\begin{proof}
Put $E=\{M:\exists k\ge0,\ \operatorname{Syr}^k(M)\le N_0\}$ and
\[
 b(z)=U_{z^\alpha}
 \{N:T_z(N)<\infty,\ \operatorname{Pass}_z(N)\in E\}.
\]
An orbit whose finite passage below $z$ belongs to $E$ has its
earlier passage below $z^\alpha$ in $E$ as well. Apply this inclusion
to $U_{z^{\alpha^2}}$, then use
Corollary~\ref{prop:Allikvere-uniform-stabilisation} to compare the two laws
at the same target $z$. Lemma~\ref{lem:Allikvere-nested-passage} gives
\[
 b(z^\alpha)\ge b(z)-C_c(\log z)^{-c}-Cz^{-\eta}
              \ge b(z)-C'_c(\log z)^{-c}                 \tag{*}
\]
for sufficiently large $z$. In particular the failure convention is
not being used to count an orbit as successful. Allikvere's final
equality with $b(z)$ in the display following \texttt{sec:assembly}
absorbs this failure mass into its $O$-term, which is legitimate.
The inequality here makes that bookkeeping explicit; it is not a
counterexample to the source argument. Tao's corresponding display in
\cite[\S3]{Tao2026v7} already uses that inequality.

Bounded $N_0$ are covered by increasing the constant. Assume henceforth
that $N_0$ is large enough that $N_0^{1/\alpha^3}$ exceeds every scale
threshold used above. If $y^\alpha\le N_0$, the assertion has witness
$k=0$ for every input. Otherwise let $J\ge1$ be least such that
\[
 v=y^{\alpha^{-J}}<N_0^{1/\alpha}.
\]
Then $N_0^{1/\alpha^2}\le v<N_0^{1/\alpha}$ and
$y=v^{\alpha^J}$. Set $z_j=v^{\alpha^{j-1}}$ for $0\le j\le J$.
Thus $z_0=v^{1/\alpha}\ge N_0^{1/\alpha^3}$,
$z_J=y^{1/\alpha}$ and $z_{j+1}=z_j^\alpha$.
At $z_0$ a finite passage lands at most $z_0<N_0$, so
$b(z_0)\ge1-Cz_0^{-\eta}$. Applying (*) successively gives
\[
 1-b(z_J)\le Cz_0^{-\eta}
   +C'_c\sum_{j=0}^{J-1}(\log z_j)^{-c}
 \le Cz_0^{-\eta}
   +\frac{C'_c(\log z_0)^{-c}}{1-\alpha^{-c}}
 \ll_c(\log N_0)^{-c}.
\]
The source of $b(z_J)$ is exactly $U_y$. Every input counted by $b(z_J)$
has an actual iterate at most $N_0$, proving the assertion.
The factor $(1-\alpha^{-c})^{-1}$ is retained in the constant; no
uniformity as $c\downarrow0$ or $\alpha\downarrow1$ is asserted.
\end{proof}

\subsection{Keeping the time budget through the same iteration}

The distribution comparison concerns passage \emph{values}, not the
joint distribution of values and passage times. To carry time, the target
set itself must contain a subsequent time-bounded orbit statement.
This is Allikvere's timed-target construction
\cite[source labels \texttt{sec:time}, \texttt{prop:timed}]{Allikvere2026v2}; the deterministic
inclusion below explains why value-distribution comparison suffices.

Put $p=4/5$ and
\[
 \Delta_0=(1-\alpha^{-3})\frac{\log N_0}{d},\qquad
 E_z^\bullet=\left\{M\in O_z:\exists k\in\mathbb N,\ 
 k\le\frac{(\log(M/N_0))_+}{d}+\Delta_0+C_1(\log z)^p,\
 \operatorname{Syr}^k(M)\le N_0\right\}.
\]
Here $(t)_+=\max(t,0)$; the nonnegative integer $k$ includes zero.
Let $B_z^\bullet$ be the event, on input $N\sim U_{z^\alpha}$, that
\[
 T_z(N)\le\frac{\log(N/z)}d+C_1(\log z)^p,\qquad
 \operatorname{Pass}_z(N)\in E_z^\bullet .
\]
Its first inequality requires finite passage.

\begin{lemma}[The timed inclusion]
\label{lem:Allikvere-timed-step}
There are constants $C_1,z_*>0$, independent of $N_0$ and $c$, for which
\[
 \mathbb P(B_z^\bullet)\ge1-O((\log z)^{-10})
       \quad (z_*\le z\le N_0),
\]
and, for $z\ge\max(z_*,N_0^{1/\alpha^3})$,
\[
 \mathbb P(B_{z^\alpha}^\bullet)
       \ge\mathbb P(B_z^\bullet)-O_c((\log z)^{-c}).
\]
\end{lemma}
\begin{proof}
The first estimate follows from Lemma~\ref{prop:Allikvere-uniform-stabilisation}: on its
concentration event the time inequality holds, and every passage
value is at most $z\le N_0$, with subsequent witness zero.

For the second estimate take the common input $N\sim U_{z^{\alpha^2}}$.
Intersect its concentration events at targets $z$ and $z^\alpha$,
and call the intersection $G$. Both are instances of
Lemma~\ref{prop:Allikvere-uniform-stabilisation}, so
$\mathbb P(G^c)\ll e^{-\kappa(\log z)^{1/5}}\ll(\log z)^{-10}$.
On $G$ both passages exist, and, with
$M=\operatorname{Pass}_{z^\alpha}(N)$ and
$M'=\operatorname{Pass}_z(N)$, the nested-passage lemma gives
\[
 t:=T_z(M)=T_z(N)-T_{z^\alpha}(N).
\]
Subtracting the two passage-time estimates and applying the drift
estimate at $T_{z^\alpha}(N)$ gives respectively
\[
 t\le\frac{\log(z^\alpha/z)}d+C(\log z)^{3/5},
 \qquad
 \log M\ge\log z^\alpha-C(\log z)^{3/5}.
\]
The drift estimate is applicable at that iterate because each passage
lies within its corresponding length-$n_0$ valuation window on $G$.
It follows that
\[
 t\le\frac{\log(M/z)}d+C'(\log z)^{3/5}.                 \tag{**}
\]
In particular this estimate concerns the same orbit, not a new
uniform sample at $M$.

Suppose now that $G$ holds and $M'\in E_z^\bullet$. If $M\le N_0$,
then $M\in E_{z^\alpha}^\bullet$ with witness zero.
If $M>N_0$ and $z\ge N_0$, concatenate the leg of length $t$ with
a witness $k_2$ for $M'$. Since $M'\le z$,
\[
 t+k_2\le\frac{\log(M/N_0)}d+\Delta_0
                      +C_1(\log z)^p+C'(\log z)^{3/5}.
\]
If $M>N_0$ and $z<N_0$, use $k_2=0$; in this case
$\log(N_0/z)/d\le\Delta_0$ by $z\ge N_0^{1/\alpha^3}$,
so (**) gives the same bound. Choose $C_1$ and $z_*$ so that, for
$z\ge z_*$,
\[
 C'(\log z)^{3/5}
       \le C_1(\alpha^p-1)(\log z)^p.
\]
Thus in both cases $M\in E_{z^\alpha}^\bullet$.
The upper estimate for $T_{z^\alpha}(N)$ also supplies the time clause
of $B_{z^\alpha}^\bullet$, after increasing $C_1$ if necessary.
We have proved the actual inclusion
\[
 G\cap\{\operatorname{Pass}_z(N)\in E_z^\bullet\}
                  \subseteq B_{z^\alpha}^\bullet.
\]
Apply the first-passage law comparison to the \emph{fixed} subset
$E_z^\bullet\subset O_z$, subtract $\mathbb P(G^c)$, and use
$B_z^\bullet\subseteq
\{\operatorname{Pass}_z(U_{z^\alpha})\in E_z^\bullet\}$.
This proves the probability inequality. No independence between
passage location and time was assumed.
\end{proof}

\begin{proposition}[A timed estimate on every power interval]
\label{prop:Allikvere-timed-block}
There is an absolute constant $C_K$ such that, with
\[
 K(N)=\frac{\log N}{d}+C_K(\log(N+2))^{4/5},
\]
for every integer $N_0\ge2$ and every sufficiently large $y$,
\[
 U_y\{N:\nexists k\in\mathbb N,\ k\le K(N),\ \operatorname{Syr}^k(N)\le N_0\}
                   \ll_c(\log N_0)^{-c}.
\]
The clock $K$ is independent of $N_0$ and $c$.
\end{proposition}
\begin{proof}
Bounded targets are absorbed into the constant as before, and
$y^\alpha\le N_0$ is trivial. Otherwise use exactly the ladder
$z_0,\ldots,z_J$ in Proposition~\ref{prop:Allikvere-untimed-ladder}.
All scales exceed $N_0^{1/\alpha^3}$, and $z_0<N_0$.
The first part of Lemma~\ref{lem:Allikvere-timed-step} starts the induction;
its second part and the same geometric error sum show
$\mathbb P(B_{z_J}^\bullet)\ge1-O_c((\log N_0)^{-c})$.
Here $z=z_J=y^{1/\alpha}$ and the input law is $U_y$.

For an input in that event, choose $k=0$ if $N\le N_0$.
Otherwise, if $z\ge N_0$, concatenate the first passage with the
$E_z^\bullet$ witness and bound
$(\log(\operatorname{Pass}_z(N)/N_0))_+\le\log(z/N_0)$.
If $z<N_0$, the passage itself is a witness, and
$\log(N_0/z)/d\le\Delta_0$. In either case an actual witness satisfies
\[
 k\le\frac{\log(N/N_0)}d+\Delta_0+2C_1(\log z)^p
   =\frac{\log N}d-\frac{\alpha^{-3}\log N_0}d
                              +2C_1(\log z)^p
   \le K(N).
\]
We used $N\ge y=z^\alpha>z$ and may take $C_K=2C_1$.
This also handles the band $N_0^{1/\alpha}\le y<N_0$ without
silently assuming that every input is already above $N_0$.
\end{proof}

\subsection{Exact coverage and the raw Collatz clock}

A single top interval would leave a proportion
$O(X^{1/\alpha-1})=O(X^{-1/1001})$ untreated. That term is not
necessary: the timed estimate holds on \emph{every} power interval,
and the intervals below the top one can be included too.

\begin{lemma}[Disjoint coverage with closed-interval estimates]
\label{lem:Allikvere-power-cover}
Let $\alpha>1$, $H\ge1$, and suppose
$t^\alpha-t>2$ for every $t\ge H^{1/\alpha}$.
Let $A$ be a set of odd integers, disjoint from $[1,H]$, such that
\[
 \#(A\cap[t,t^\alpha])
       \le\varepsilon\,\#((2\mathbb N+1)\cap[t,t^\alpha])
       \quad(t\ge H^{1/\alpha}).
\]
Then for every real $X\ge1$,
\[
 \#(A\cap[1,X])\le
           2\varepsilon\,\#((2\mathbb N+1)\cap[1,X]).
\]
\end{lemma}
\begin{proof}
For $X\le H$ the left side vanishes. Otherwise put
$t_j=X^{\alpha^{-j}}$ and let $J$ be least such that $t_J\le H$.
For $0\le j<J$, $t_{j+1}>H^{1/\alpha}$ and
$t_j=t_{j+1}^\alpha$.
The intervals $(t_{j+1},t_j]$ are pairwise disjoint and cover
$(t_J,X]$. Each has length greater than two, hence contains at
least one odd integer. Passing from this half-open interval to
$[t_{j+1},t_j]$ adds at most its lower endpoint, hence at most
one integer. Its closed-interval odd count is therefore at most
twice its half-open odd count. Sum the given estimate over the
closed intervals and then the disjoint half-open counts.
The remaining segment $[1,t_J]$ is contained in $[1,H]$ and
contains no element of $A$. This proves the assertion, with
all real or integer endpoint coincidences included.
\end{proof}

\begin{corollary}[Timed fixed-target count, without a cutoff error]
\label{cor:Allikvere-timed-count}
For every $0<c<1/17.232$ there is $C_c$ such that for every real
$X\ge2$ and every integer $N_0\ge2$,
\[
 \frac1X\#\{N\le X:N\ {\rm odd},\
       \nexists k\in\mathbb N,\ k\le K(N),\ \operatorname{Syr}^k(N)\le N_0\}
          \le C_c(\log N_0)^{-c}.                      \tag{***}
\]
In particular $c=1/18$ is admissible. Dropping the time restriction
gives the same estimate for orbits that never attain $N_0$.
\end{corollary}
\begin{proof}
For sufficiently large $N_0$, apply Lemma~\ref{lem:Allikvere-power-cover}
with $H=N_0$ to the failure set in Proposition~\ref{prop:Allikvere-timed-block}.
It is disjoint from $[1,N_0]$ because time zero is allowed.
The elementary function $t^\alpha-t$ is increasing for $t\ge1$
and tends to infinity, so the width hypothesis holds after one
fixed lower bound on $N_0$. All source intervals are then within
the range of the proposition. Finally
$2\#((2\mathbb N+1)\cap[1,X])\le X+1\le3X/2$ for $X\ge2$.
For targets below that fixed lower bound, enlarge $C_c$ so that
$C_c(\log N_0)^{-c}\ge1$. This covers every target and cutoff.
\end{proof}

This is a consequence of Allikvere's timed iteration with the stronger
kernel estimate proved above, and of the scale-covering mechanism already
used by Tao in \cite[\S3]{Tao2026v7}. Relative to Allikvere's stated
Theorem~\texttt{thm:quant}, it replaces the displayed exponent $1/29$
by every $c<1/17.232$ and removes the extra $X^{-1/2000}$ term.
The removal is a counting consequence, not a new mixing theorem.

Define the unaccelerated map $C$ by $C(n)=3n+1$ for odd $n$
and $C(n)=n/2$ for even $n$. These are the individual steps counted
in the following statement.

\begin{corollary}[One threshold-independent clock]
\label{cor:Allikvere-universal-clock}
There is a constant $C_R$ for which the clock
\[
 R(N)=\left(\frac3d+\frac1{\log2}\right)\log N
                         +C_R(\log(N+2))^{4/5}
\]
satisfies, for every $X\ge2$, $N_0\ge2$ and $0<c<1/17.232$,
\[
 \frac1X\#\{N\le X:\nexists m\in\mathbb N,\ m\le R(N),\
                       C^m(N)\le N_0\}
              \ll_c(\log N_0)^{-c}.
\]
Consequently, for each function $f:\mathbb N_{\ge1}\to\mathbb R$
with $f(N)\to+\infty$,
\[
 \#\{N\le X:\nexists m\le R(N),\ C^m(N)<f(N)\}=o(X).
\]
For positive odd inputs the same conclusion holds with $\operatorname{Syr}$ and
$K$ in place of $C$ and $R$. The clocks are independent of $f$;
the exceptional set is permitted to depend on $f$.
\end{corollary}
\begin{proof}
For odd $M$ write $M_i=\operatorname{Syr}^i(M)$ and
$s_k=\sum_{i=0}^{k-1}v_2(3M_i+1)$. Concatenating the actual
odd step and its subsequent halvings proves
$C^{k+s_k}(M)=M_k$. Moreover
\[
 2^{s_k}M_k=\prod_{i=0}^{k-1}(3+1/M_i)\,M
              \le4^kM,
 \qquad s_k\le2k+\log_2 M.
\]
The empty product covers $k=0$. For an arbitrary positive integer
$N=2^aM$ with $M$ odd, an odd-map witness $k\le K(M)$ therefore gives
a raw-map witness at
\[
 m=a+k+s_k\le3k+a+\log_2 M=3k+\log_2 N
       \le R(N)
\]
with $C_R=3C_K$. The exact coordinate map
$M\mapsto2^aM$ identifies odd integers $M\le X/2^a$ with inputs
$N\le X$ of valuation exactly $a$. Hence the raw failure count is
at most the sum of the odd failure counts at cutoffs $X/2^a$.
By (***) that sum is at most
\[
 C_cX(\log N_0)^{-c}\sum_{a\ge0}2^{-a}
          =2C_cX(\log N_0)^{-c}.
\]
Cutoffs below two contribute nothing, since their only possible
positive odd input is $1\le N_0$. This proves the uniform fixed-target
estimate without a separate tail-limit argument.

For a diverging $f$, put
$b(X)=\inf_{n\ge\sqrt X}f(n)$. For large $X$ this is finite and
exceeds three, and $b(X)\to\infty$. Choose the integer
$N_0(X)=\lceil b(X)\rceil-1$, so that
$2\le N_0(X)<b(X)$ and $N_0(X)\to\infty$.
Every input $N\ge\sqrt X$ that reaches $N_0(X)$ reaches a value
strictly less than $f(N)$. The failure proportion is therefore at
most
\[
 X^{-1/2}+O_c((\log N_0(X))^{-c})=o(1).
\]
The identical argument uses (***) for odd inputs. In particular,
discarding the time bound recovers Allikvere's natural-density
conclusion. This is a written reconstruction from the cited analytic
inputs, not an inference from the finite trajectory tests.
\end{proof}

The exact Gaussian coordinate calculation above also evaluates Mazur's
original raw tube mass. Section~\ref{sec:Mazur-mass-interior} uses precisely
half this envelope, retains its different strict cutoff and both tails,
and proves the physical interior profile rate
$O_C(N^{-1/17.232}(\log N)^{1/2-1/17.232})$.
The comparison retains the exact Robbins remainder and floor correction;
it is not an identification of the two source kernels or sampling laws.

Section~\ref{sec:Mazur-genuine-assembly} now completes the comparison
with actual first-passage laws on Mazur's two original closed source
blocks. Both harmonic and uniform sampling have common-center error
$O((\log B)^{-1/17.232}(\log\log B)^{1/2-1/17.232})$.
The full-prefix bound makes the original tube family usable at $C=1/2$;
the source's final assembly takes $C=40$. The endpoint kernels have not
been identified across the two authors: the comparison is between
their actual passage-law statements, and the maps and exceptional
masses are retained. Section~\ref{sec:Mazur-real-timed} now reconstructs
the subsequent real-source and timed iteration, preserving the critical
rate and yielding Syracuse and ordinary clocks $100\log_2 n$ and
$301\log_2 n$ within Mazur's schedules. These refine that source's
constants, not the much smaller leading Syracuse clock proved above.

\paragraph{Comparison with the reconstructed timeout argument.}
Section~\ref{sec:Shaik-timeout} gives a different, polylogarithmic target
with same-witness height control. It has shortcut, ordinary and
odd-return leading clocks \(2,3,1\) divided by \(\log(4/3)\), with
additive error \(O(\sqrt{\log n\,\log\log n})\).
The fixed-target/unrestricted diverging-target statement above is
retained with its own exceptional estimate and clock; the two
conclusions have not been identified.
